\section{Introduction}\label{sec:introduction}

Derivative values are conditional on parameters and state variables that are not observed without error. In the Black--Scholes--Merton benchmark, volatility is often inserted into the pricing map as though it were known, even when it has been estimated from a finite return history. A Bayesian treatment replaces that point estimate with a posterior distribution and can propagate posterior uncertainty through the derivative-pricing map. This idea is not new to option pricing: Bayesian learning has been used to explain implied-volatility dynamics and improve out-of-sample option forecasts \cite{guidolin2003}, risk-neutral predictive densities have been constructed while integrating parameter uncertainty \cite{romboutsstentoft2014}, and Bayesian calibration has been used to propagate uncertainty in calibrated model parameters to prices of exotic derivatives \cite{guptareisinger2014}. The question in this paper is therefore narrower: \emph{when does posterior integration materially change a point price, and how important is that effect relative to changing the volatility or pricing specification itself?}

Average-price commodity options provide a useful setting for that comparison. Their payoff depends on a path of underlying prices, so the curvature of the pricing map can change materially with moneyness, maturity, and the fraction of the averaging window already fixed. The numerical literature has long exploited geometric-average structure to price arithmetic-average claims efficiently \cite{kemnavorst1990,curran1994}. Commodity-specific models add mean reversion, convenience yield, stochastic interest rates, and jumps \cite{schwartz1997,ewaldwu2023}. Most importantly for the present application, \cite{shirayatakahashi2011} study WTI average-price options under Heston and extended $\lambda$-SABR dynamics, explicitly calibrating richer risk-neutral models to more liquid American vanilla options on WTI futures before valuing the average option. \cite{ganwangyang2020} provide a distinct model-free machine-learning application using real WTI Average Price Option data. These contributions mean that neither WTI APO pricing nor the use of richer calibrated risk-neutral dynamics is itself the novelty here.

The empirical distinction between parameter uncertainty and model specification is also central in the broader option-pricing literature. Richer stochastic-volatility and jump models can reduce pricing errors, but additional complexity does not improve every performance criterion \cite{bakshicaochen1997,cumminsesposito2025}. Flexible implied-volatility functions can fit the contemporaneous cross-section extremely well and still deteriorate out of sample, which makes forward validation essential \cite{dumasflemingwhaley1998}. In crude oil specifically, option-implied volatility measures contain forward-looking information beyond historical return models \cite{gildertsiaras2020}, commodity variance risk premia imply an economically meaningful distinction between physical and risk-neutral variation \cite{teeting2017}, and volatility-of-volatility risk is itself priced in oil-option returns \cite{roh2021}. These results motivate treating an option-implied scalar volatility as an \emph{effective} or model-equivalent risk-neutral state rather than as a structural diffusion coefficient.

This paper studies that comparison using CME WTI Average Price Options (APOs). The contract pays against the arithmetic average of daily first-nearby WTI futures settlements over a calendar month. Consequently, a contract-consistent valuation exercise must account for the first-nearby roll convention and partial fixing: once the averaging month has begun, some components of the final average are known while the remaining components are stochastic. Treating the option as an ordinary European claim on a single fixed-maturity futures contract would change the object being valued.

The central methodological principle is a strict separation between inference under the physical measure and valuation under the risk-neutral measure. Historical returns are used to infer a data-generating process under $\mathbb P$. Posterior uncertainty in volatility is then mapped into option values under $\mathbb Q$. The physical drift is not inserted into the risk-neutral pricing dynamics. This distinction follows the no-arbitrage separation between physical inference and risk-neutral valuation \cite{blackscholes1973,merton1973}. Under the maintained constant-volatility benchmark, the same diffusion coefficient can coherently be used under both measures; whether that historical volatility is empirically adequate for the option market is a separate specification question.

Let $C^{Q}(\sigma)$ denote the risk-neutral value of an arithmetic-average option conditional on volatility $\sigma$. The posterior-integrated rule is
\begin{equation}
 C_{PI}=\mathbb E\!\left[C^{Q}(\sigma)\mid\mathcal D\right],
 \label{eq:intro_fb}
\end{equation}
whereas the posterior-mean plug-in rule is
\begin{equation}
 C_{PM}=C^{Q}\!\left(\mathbb E[\sigma\mid\mathcal D]\right).
 \label{eq:intro_pm}
\end{equation}
The two need not coincide because the option-pricing map is nonlinear. A second-order expansion around $\bar\sigma=\mathbb E[\sigma\mid\mathcal D]$ gives
\begin{equation}
 C_{PI}-C_{PM}
 \approx
 \frac{1}{2}C^{Q\prime\prime}(\bar\sigma)
 \operatorname{Var}(\sigma\mid\mathcal D),
 \label{eq:intro_taylor}
\end{equation}
which isolates two forces: posterior dispersion and curvature of the pricing function. Parameter uncertainty should matter most where the posterior is diffuse and option value is highly curved in volatility.

The paper develops four layers of evidence. First, repeated-sampling experiments and a 175,000-history mechanism map provide controlled environments in which the true volatility is known. The mechanism map varies historical sample size from 21 to 1,260 observations, true volatility from 0.10 to 0.80, moneyness from 0.60 to 1.50, maturity from 21 to 252 days, and the fraction of the averaging window already fixed. This makes it possible to identify where the approximation in Equation~\eqref{eq:intro_taylor} succeeds and where posterior integration can become economically material.

Second, the paper develops a contract-consistent empirical application using historical WTI APO settlements obtained through Barchart. The downloaded database contains 213 unique call and put contracts and 11,179 option-date observations across expiries from September 2026 through June 2029. The October-2026 baseline contains 310 retained option-date observations across 12 valuation dates from August 25 through September 10, 2026. Additional September and October panels provide 402 contract-date observations for implied-volatility and heavy-tail robustness analyses. Barchart \texttt{Latest} is treated as the end-of-day settlement supplied in the source histories; it is not asserted to be an independently verified official CME settlement.

Third, the paper distinguishes posterior integration from volatility specification using the APO market itself. It inverts the maintained pricing map to obtain Curran-equivalent APO-implied volatilities and then asks whether volatility smiles estimated only from earlier dates improve subsequent APO valuation. This strict forward-in-time design spans seven expiries and prevents same-day target prices from entering the fitted smile, but the volatility information still comes from the same option family.

Fourth, the paper breaks that remaining within-family link using standard monthly WTI futures options. Official prior-date LO option settlements and matching CL futures settlements are inverted through an American futures-option model to obtain an independent risk-neutral volatility surface. That surface is mapped to the APO fixing horizon and used to price October-2026 APOs without using APO prices to construct the external state. This cross-instrument experiment asks whether the pricing gain attributed to risk-neutral volatility information survives when the source and target option families are separated.

The main results separate the two effects sharply. Within the October baseline, posterior-integrated and posterior-mean pricing are almost identical: pooled RMSE is 0.5177 versus 0.5183 across all 310 observations and 0.1864 versus 0.1869 in the six positive-volume contract-date observations. A high-precision benchmark using 43.56 billion pseudo-random path--volatility evaluations and an independent 16.24 billion randomized Sobol evaluations confirms that the observed PI--PM differences of roughly $10^{-3}$ are not Monte Carlo noise. The mechanism map nevertheless finds states in which posterior integration is material: the largest observed absolute gap is 0.1384, concentrated in short historical samples, high volatility, long horizons, extreme moneyness, and limited realized fixing.

The larger empirical effect is volatility specification. The extended APO panel reveals a pronounced maturity structure in the effective volatility required by the maintained pricing map: median date-level common $\sigma_Q$ is approximately 0.448--0.457 for the September--October 2026 expiries, 0.361 for March 2027, 0.300 for September 2027, 0.265 for March 2028, and 0.242 for September 2028, while the historical posterior mean remains close to 0.416. Strict forward-in-time validation uses only implied-volatility observations dated before each target date. Across 2,381 holdout option-date observations spanning seven expiries, the historical-volatility posterior-integrated baseline has MAE 2.6187 and RMSE 3.4090, compared with MAE 0.1088 and RMSE 0.1725 for the expanding prior-date smile and MAE 0.1075 and RMSE 0.1594 for the previous-day smile.

The independent vanilla-WTI experiment shows that this result is not confined to information extracted from the APO family. After expanding the observed LO strike range to 70--120, 253 of 280 matched October-2026 contract-dates lie inside the prior LO moneyness support. On that common-support sample, historical-volatility posterior integration has MAE 0.4464 and RMSE 0.5562, compared with 0.1559 and 0.2599 for the external previous-day LO surface. The corresponding prior-date APO smile has MAE 0.1676 and RMSE 0.2686. Valuation-date cluster bootstraps show a large and stable improvement of external LO over historical volatility; differences between the external LO and prior-date APO surfaces are small, supporting comparability rather than a claim of universal vanilla-option superiority.

The paper therefore contributes a conditional rather than universal result. It does not propose Bayesian option pricing, WTI APO pricing, or stochastic-volatility calibration as new ideas. Instead, it isolates a specific within-model mechanism: the incremental point-price effect of posterior integration is approximately governed by posterior dispersion times pricing curvature. The mechanism map identifies the states in which that correction becomes economically material, while the WTI application places the correction beside the empirically much larger volatility-specification margin. In the observed WTI APO sample, posterior integration is second-order relative to the difference between a long-window historical volatility estimate and an option-informed risk-neutral volatility state. The independent LO experiment shows that this ordering does not require calibrating and evaluating within the same APO family. In weak-information, high-curvature states, however, the synthetic mechanism map shows that the ordering can reverse. This comparison provides a benchmark for deciding when Bayesian propagation is worth carrying before adding further model complexity.

The remainder of the paper is organized as follows. Section~\ref{sec:contract} describes the WTI Average Price Option and the empirical panel. Section~\ref{sec:methodology} presents physical-measure inference, risk-neutral valuation, and the competing pricing rules. Section~\ref{sec:design} describes the synthetic, empirical, and forward-validation designs. Section~\ref{sec:results} reports the mechanism and market evidence. Section~\ref{sec:robustness} evaluates numerical accuracy, posterior calibration, heavy-tailed returns, prior sensitivity, and liquidity dependence. Section~\ref{sec:conclusion} concludes.
